\documentclass[12pt,letterpaper]{article}

\usepackage[margin=1in]{geometry}
\usepackage[T1]{fontenc}
\usepackage{amsmath}
\usepackage{amssymb}
\usepackage{newtxtext,newtxmath}
\usepackage{setspace}
\usepackage{caption}
\usepackage{graphicx}
\usepackage{titlesec}
\usepackage[hidelinks]{hyperref}

% Match the blueprint's body-size, bold, unnumbered section headings.
\titleformat{\section}{\normalfont\normalsize\bfseries}{}{0pt}{}
\titlespacing*{\section}{0pt}{\baselineskip}{0pt}
\setlength{\parindent}{0pt}
\setlength{\parskip}{0pt}

\begin{document}

\hypersetup{pageanchor=false}
\begin{titlepage}
    \centering
    \singlespacing
    \vspace*{1.5in}

    {\bfseries Does higher education complement or substitute for family background?\par}
    \vspace{\baselineskip}
    Econ 4331W Development Economics\par
    November/11/2026\par
    Ellie Hou\par

    \vfill
\end{titlepage}
\hypersetup{pageanchor=true}

\doublespacing
\raggedright

\section*{INTRODUCTION}
Does completing a four-year college degree give students from a disadvantaged background a
larger earnings premium? 


This question helps inform a policy trade-off: to help students from
disadvantaged backgrounds thrive, do we simply lower tuition, or is a reform of the whole
education system needed? With data from the National Longitudinal Survey of Youth 1997,
this paper examines the difference in the college earnings premium by family background,
controlling for pre-college academic preparation, measured by ASVAB scores, and other
demographic characteristics.

I define financial disadvantage using baseline household income, which measures
resources available before college. I distinguish completed degree levels so that
the bachelor's premium is measured against high-school diploma holders rather
than a mixture of high-school dropouts, diploma holders, and associate graduates.


\section*{LITERATURE REVIEW}

Mincer (1974) links earnings to schooling and work experience, treating these as human
capital investments in earning capacity. The framework provides a foundation for answering
my main research question by measuring how schooling can affect earning capacity.
Furthermore, Becker and Tomes (1986) asked a broader question: how do parents influence
their children’s human capital? In their model, parents prefer to optimize the amount of
investment in human capital before starting to save for voluntary bequests.
These questions also connect to the literature on wealth concentration and inequality. De
Nardi (2015) reviews basic Bewley-style models, in which households save to protect
consumption against earnings risk, alongside models incorporating bequests and
intergenerational human-capital transmission. The model helps explain the role of family investment in
differences in earning capacity. In this paper, I will focus on a narrower question: “Is college
completion associated with narrowing income disparities among different family-background
groups?”
More directly, Hendricks, Koreshkova, and Leukhina (2026) study the potential effect of
widening access to selective public universities on earning capacity. Their results show that
expanding enrollment capacity can increase the population’s aggregate earning capacity,
partly by admitting academically well-prepared students who were previously excluded by
admission constraints. My analysis, on the other hand, examines whether the adjusted
college earnings premium varies with family background.

Chetty, Deming, and Friedman (2026) study admission to highly selective private
colleges. Their need-affirmative admissions exercise uses the bottom quintile of
parental income as a policy threshold, following Chetty et al. (2020). I adapt
that income threshold to define a lower-income group in NLSY97. Their causal
analysis concerns college selectivity; my analysis instead estimates conditional
earnings associations across completed degree levels.

\section*{HYPOTHESIS}
My hypothesis is that students from lower-income families have a higher adjusted college
earnings premium. Specifically, I test whether the bachelor's-only versus high-school
premium is larger for the lower-income group, $\delta_{\mathrm{BA}}>0$, in the model
described in the methodology section.
The proposed mechanism behind my hypothesis is that college education can substitute
for family background and thus narrow differences in human capital among students from
different backgrounds.

% TODO: Support the proposed mechanism for the secondary hypothesis with literature.
My secondary hypothesis is that students from lower-income families are less likely to
complete a bachelor's degree. The reasoning behind this hypothesis is that, if schools
value outcomes associated with future earnings, colleges are incentivized to invest more
in students with higher demonstrated potential, and that potential can logically be
highly correlated with students' backgrounds.

\section*{DATA}
\textbf{Data source and sample.}
For this paper, I used the data from NLSY97 (the National Longitudinal Survey of Youth)
to construct variables. The survey includes a representative sample of Americans born
between 1980 and 1984. The source variables I included are household income in 1996,
parent-reported household net worth in 1997, wage income, and highest degree earned in 2015.

The revised balanced panel includes 2,849 people with usable measures in three
earnings years, giving 8,547 person-year observations. All respondents have a positive
official custom weight for this selected panel. I retain reported zero earnings and
recover a skipped wage amount as zero only when the income-receipt screening
questions explicitly establish no wage income. Unresolved skips and other missing
amounts remain excluded. This recovers 1,402 zero wage amounts and adds 726 people
relative to the preliminary sample. Genuine negative baseline income and net worth
are retained. For comparability, the sample continues to require observed baseline
net worth, although wealth no longer defines the income groups.

\textbf{Dependent variable $Y_{it}$.}
$Y_{it}$ is the annual wage and salary income of the surveyed person in 2016, 2018,
and 2020, measured in 2025 US dollars. These are the earnings years available in the
extract, collected in the 2017, 2019, and 2021 surveys.
Annual earnings reflect both wage rates and time worked; they are not hourly wage
rates. Zero wage income need not imply nonemployment or zero self-employment income.

\textbf{Education categories $D_{ik}$ and lower-income background $L_i$.}
I classify highest degree by the 2015 interview into five mutually exclusive
categories: no high-school diploma or GED only; high-school diploma; associate
degree; bachelor's degree only; and graduate or professional degree. High-school
diploma is the reference category. This separates associate, bachelor's, and
graduate credentials, while retaining pooled categories for no diploma/GED and
graduate/professional degrees. High-school diploma holders can include people who
attended college without completing a postsecondary degree. Degree status stays
fixed across all three earnings years.

$L_i=1$ if 1996 household income is at or below its weighted 20th-percentile cutoff,
\$21,360 in nominal 1996 dollars. I calculate this cutoff within the selected panel using
\href{https://www.nlsinfo.org/content/weights/nlsy97}{custom NLSY97 weights}, counting
each person once. The cutoff is the smallest observed income at which cumulative
weight reaches 20\%; all ties at the cutoff are included. The lower-income group
contains 765 people and represents 20.04\% of weighted respondents. The comparison
group consists of all incomes above the cutoff, not exclusively wealthy families.
Wealth does not enter this indicator.

This is an adaptation of the bottom-quintile threshold in Chetty, Deming, and
Friedman (2026, Section VI.A and Table V), not a replication of their income measure.
They average parental household income over child ages 12--17 and construct
national birth-cohort ranks; I use one year's household income and pool the
1980--84 birth cohorts within the selected sample. I also retain one respondent
with genuine negative baseline income, whereas their analysis excludes negative
parental incomes. The cutoff therefore describes relative income in this analytic
sample and is not an estimate of the national bottom-quintile threshold.

\textbf{Control variables $X_{it}$.}
I control for pre-college academic achievement using the NLSY97-created ASVAB
math--verbal percentile score, which is normalized within three-month birth groups.
In addition, I control for sex, age at the end of the earnings year, and race/ethnicity.
% Verified in the archived codebook: R9829600 is the age-normalized percentile;
% dividing by 1,000 restores its three implied decimal places.

% TODO: Choose the summary-table presentation after revising the sample.
\textit{[Insert your summary statistics table here.]}

\section*{METHODOLOGY}
I convert nominal dollars into constant 2025 dollars using annual average CPI-U,
multiplying each value by $CPI_{2025}/CPI_t$.\footnote{The supplied CPIAUCNS series
has no October 2025 observation, consistent with
\href{https://www.bls.gov/news.release/archives/cpi_12182025.htm}{BLS documentation}.
I use the mean of the 11 available months, 321.943.}

The primary model is
\[
\begin{aligned}
    Y_{it}={}&\alpha+\gamma L_i+\sum_{k\in\mathcal K}\beta_k D_{ik}
              +\sum_{k\in\mathcal K}\delta_k(D_{ik}\times L_i)\\
             &+X'_{it}\theta+\lambda_t+\epsilon_{it},
\end{aligned}
\]
where $\mathcal K$ contains the four education categories other than high-school
diploma, and $D_{ik}$ indicates membership in category $k$. $\beta_k$ is the adjusted
earnings difference between category $k$ and high-school diploma holders among
respondents above the income cutoff; $\beta_k+\delta_k$ is the corresponding
difference for the lower-income group. Thus $\delta_{\mathrm{BA}}$ compares the
bachelor's-only premium across income groups. $\lambda_t$ controls for earnings-year
effects, and $X_{it}$ includes ASVAB percentile, sex, age, and race/ethnicity.

I estimate weighted pooled OLS with year effects, applying the same custom person
weight in each year, and cluster standard errors by person. Confidence intervals
for $\beta_k+\delta_k$ use the full coefficient covariance. I also report a joint
test that all four interactions are zero. These are adjusted associations, not
causal difference-in-differences estimates: the model does not identify an
exogenous change in education or establish parallel trends.

The income-only 20\% definition was chosen from the literature before fitting this
revised model. It is not preregistered: the preliminary report explored several
joint income-and-wealth cutoffs and selected 60\% by its smallest interaction
$p$-value. The new citation does not retrospectively justify that earlier choice.

\section*{RESULTS}
% Numerical prose updated from outputs/regression; presentation remains for revision.
The adjusted bachelor's-only versus high-school earnings premium is \$12,860 for
the lower-income group and \$31,136 for respondents above the cutoff, in 2025 dollars.
The estimated difference is $\hat\delta_{\mathrm{BA}}=-\$18,276$
(clustered standard error \$6,828; $p=0.0075$; 95\% confidence interval:
$-\$31,664$ to $-\$4,889$). This estimate runs against my hypothesis of a larger
bachelor's premium for lower-income students. The joint test of all four education
interactions yields $F(4,2848)=2.008$ ($p=0.091$), which does not reject their joint
equality to zero at the 5\% level. These tests address different hypotheses.
The lower-income group contains only 25 graduate/professional degree holders,
so its graduate-premium estimate is imprecise.

% TODO: Choose the figure and regression-table presentation later.
\textit{[Insert your figure and its interpretation here.]}
\textit{[Insert your regression results table here.]}

Using the custom weights, 14.3\% of the lower-income group completed a bachelor's
degree or higher, compared with 42.3\% of the other group. This descriptive comparison
is consistent with my secondary hypothesis. These results describe weighted
associations within the selected balanced panel. The pointwise tests do not adjust
for multiple comparisons, estimated-cutoff uncertainty, or the full survey design.
Missing data, topcoded earnings, and unobserved differences
limit the interpretation; the analysis does not establish a causal effect of college
completion on wealth inequality.

\section*{EXTENSIONS}
The smaller bachelor's premium for lower-income respondents is consistent with
complementarity between education and family resources, but this observational
comparison does not establish that mechanism. Selection into degree completion,
differences in time worked, and the measurement of family resources could also
contribute. Future checks can examine baseline wealth separately and assess
sensitivity to the single-year income measure and sample restrictions.

% TODO: Substantiate these channels. College selectivity is not defined for
% never-attenders and may be a pathway, so adding it as a control changes the question.
One additional theory for future research is that, even if students from disadvantaged
backgrounds indeed have a lower college premium, this might be because they tend to
attend less selective schools (affected by cost, nepotism, and unobserved academic
preparation). Hence, further research with additional control variables such as the
selectivity of the college attended can be helpful.

\section*{REFERENCES}
\begin{list}{}{\setlength{\leftmargin}{0.3in}
              \setlength{\itemindent}{-0.3in}
              \setlength{\itemsep}{0.6\baselineskip}
              \setlength{\parsep}{0pt}}
\item Becker, Gary S., and Nigel Tomes. 1986. ``Human Capital and the Rise and Fall of Families.'' \textit{Journal of Labor Economics} 4 (3, Part 2): S1--S39.
\item Chetty, Raj, David J. Deming, and John N. Friedman. 2026. ``Diversifying Society's Leaders? The Determinants and Causal Effects of Admission to Highly Selective Private Colleges.'' \textit{Quarterly Journal of Economics} 141 (1): 51--145. \href{https://doi.org/10.1093/qje/qjaf050}{doi:10.1093/qje/qjaf050}.
\item Chetty, Raj, John N. Friedman, Emmanuel Saez, Nicholas Turner, and Danny Yagan. 2020. ``Income Segregation and Intergenerational Mobility Across Colleges in the United States.'' \textit{Quarterly Journal of Economics} 135 (3): 1567--1633. \href{https://doi.org/10.1093/qje/qjaa005}{doi:10.1093/qje/qjaa005}.
\item De Nardi, Mariacristina. 2015. ``Quantitative Models of Wealth Inequality: A Survey.'' NBER Working Paper 21106.
\item Hendricks, L., T. Koreshkova, and O. Leukhina. 2026. ``Expanding Access to Selective Colleges.'' Federal Reserve Bank of St. Louis Working Paper 2026-005.
\item Mincer, Jacob A. 1974. \textit{Schooling, Experience, and Earnings}. National Bureau of Economic Research.
\end{list}
\end{document}
